\subsection{Coupling all Morita action matrices}\label{sec:full-morita-gram}
The two-sided module test can be strengthened by retaining the complete
associativity pairing, rather than its restrictions to one module simple at
a time. The strengthening uses Frobenius--Perron positivity only to recognize
module categories; it makes no pivotal or spherical assumption.

\begin{theorem}[Full Morita action Gram identity]\label{thm:full-morita-gram}
Let ${}_{\mathcal C}\mathcal M_{\mathcal D}$ be an invertible bimodule
category, with simple objects $m_0,\ldots,m_{s-1}$. Write
\[
 L_i(a,b)=[X_i\otimes m_a:m_b],\qquad
 R_j(a,b)=[m_a\otimes Y_j:m_b]
\]
for its left and right action matrices. Then
\begin{equation}\label{eq:full-morita-gram}
 \sum_j R_j(a,d)R_j(b,c)
       =\sum_i L_i(b,a)L_i(c,d)
       \qquad(0\leq a,b,c,d<s).
\end{equation}
Consequently the matrix
\[
 G_{(a,d),(b,c)}=\sum_i L_i(b,a)L_i(c,d)
\]
has a nonnegative integral Gram factorization whose columns are the
vectorized matrices $R_j$. Each such column belongs to
$\operatorname{im}_{\mathbb Q}G$, and its coordinate indexed by $(a,d)$ is
at most $\lfloor\sqrt{G_{(a,d),(a,d)}}\rfloor$.
The corresponding matrices commute with every $L_i$, include the identity,
and are permuted by transposition.
\end{theorem}
\begin{proof}
Use duals in the Morita bicategory. The multiplicity of $Y_j$ in
$m_b^*m_c$ is $R_j(b,c)$, by Frobenius reciprocity. Thus the multiplicity
of $m_d$ in $m_a(m_b^*m_c)$ is the left-hand side of
\eqref{eq:full-morita-gram}. The multiplicity of $X_i$ in $m_am_b^*$ is
$L_i(b,a)$, so associativity identifies the same multiplicity with its
right-hand side. If $W$ is the matrix of vectorized right actions, this
is $G=WW^{\mathsf T}$. Hence
$\operatorname{im}_{\mathbb R}W=\operatorname{im}_{\mathbb R}G$; since all
matrices are rational, each column of $W$ belongs to the stated rational
image. A diagonal entry of $G$ is the sum of the squares of the
corresponding entries of $W$, which proves the bound. The commuting
left and right actions, the tensor unit of $\mathcal D$, and
$R_{j^*}=R_j^{\mathsf T}$ prove the remaining assertions.
\end{proof}

Suppose now that $A=\mathbf1\oplus g$ is the split involution algebra and
$m_0=A$ in $\mathcal C_A$. The existence argument for this algebra and the
exhaustive module-action construction are those of the preceding module
criterion. Since the internal end of $m_0$ is $A$, setting $c=d=0$ in
\eqref{eq:full-morita-gram} gives
\begin{equation}\label{eq:morita-first-columns}
 \sum_j R_j(\,\cdot\,,0)R_j(\,\cdot\,,0)^{\mathsf T}=I+L_g.
\end{equation}
Here $L_g$ is an involutive permutation matrix. Its right-hand side has
a unique nonnegative integral Gram factorization with nonzero columns,
up to permutation of columns: a fixed point contributes two copies of its unit vector, and
a transposition contributes the sum of its two unit vectors. Indeed,
zero off-diagonal entries prevent columns from meeting different
orbits. On a fixed point the squared entries sum to two; on a
transposition the diagonal entries and the off-diagonal entry are all
one. These observations prove the assertion orbit by orbit.
Transposition supplies the same finite list for the possible first rows
of each right action matrix.

The implementation intersects the rational image of $G$ with these
prescribed first rows and columns. Choosing free coordinates among the
actual matrix entries, the diagonal bounds in
Theorem~\ref{thm:full-morita-gram} give a finite integer box. Every
candidate is checked for nonnegativity, the entrywise necessary bounds
$vv^{\mathsf T}\leq G$, and commutation with the left actions. If the
candidates are $v_1,\ldots,v_t$, the unknown multiplicities must satisfy
\[
 \sum_{h=1}^t z_hv_hv_h^{\mathsf T}
       =G-\operatorname{vec}(I)\operatorname{vec}(I)^{\mathsf T},
 \qquad z_h\in\mathbb Z_{\geq0},
\]
together with the first-column multiplicities from
\eqref{eq:morita-first-columns}, after removing the unit action.
An exact Farkas certificate for the relaxation $z_h\geq0$ excludes that
module-action branch. As before, exclusion of the fusion ring requires
exhausting every eligible branch. The program
\filename{software/full_morita_gram.py} reconstructs the finite domains
and checks separating vectors using integer arithmetic.

\begin{proposition}[The involution dual of the first residual]
\label{prop:first-residual-morita-self-ring}
If $\mathcal C$ categorifies $\RR{3}{215}$ and $g$ is its nontrivial
invertible simple, then $\mathbf1\oplus g$ admits a connected separable
algebra structure and
\[
 K_0\bigl({}_A\mathcal C_A\bigr)\cong K_0(\mathcal C)
 \quad\text{as based rings},\qquad A=\mathbf1\oplus g.
\]
This assertion does not identify the two fusion categories up to tensor
equivalence.
\end{proposition}
\begin{proof}
The involution fixes two simples, so the previously established split
involution argument applies. Exhaustive reconstruction gives sixteen
module-action branches. Two violate the rank bound
$\operatorname{rank}G\leq\operatorname{rank}K_0({}_A\mathcal C_A)$,
and twelve have exact Farkas certificates for the full Gram system.
Each of the remaining two has Gram rank six and six dual simples.
Therefore its six right action matrices are linearly independent; in
particular they cannot repeat. Exhaustive selection from the finite
candidate lists gives exactly one Gram frame for each branch, up to
permuting its columns. Multiplying those six matrices reconstructs a
nonnegative integral fusion tensor, and a recorded based permutation
identifies all its coefficients with $\RR{3}{215}$.

The complete sixteen-branch record is
\filename{research/new_frontier/frontier_full_morita.json}.
The script
\filename{research/new_frontier/verify_frontier_full_morita.py}
reconstructs the module actions and candidate domains, checks all rank
and Farkas obstructions, proves uniqueness of each surviving frame by
finite exact enumeration, and verifies every matrix product and the
based-ring identification. Verification does not call an optimizer.
\end{proof}

The two surviving frames are based bimodule data, not associators or
categorifications. In particular the proposition does not supply an
algebra $\mathbf1\oplus X$ for a noninvertible simple $X$.
The uniform full Gram scan covers all $47$ residual inputs through
multiplicity eight: $21$ are eligible for the implemented involution
search and $26$ are ineligible. It yields no additional whole-ring
exclusion. All $54$ supplied positive models pass the control scan.
These outcomes are retained separately from the exclusion ledger.
